% 07_experiments.tex
\section{Experimental Validation}
\label{sec:experiments}

We validate MCL on a curated benchmark of 173 Wikipedia concepts
covering notable figures, historical events, scientific phenomena,
geographical entities, technology, and popular culture. Each concept
yields a prompt of the form \emph{``Explain [concept] in a
comprehensive way.''} Generation uses greedy decoding for
watermarked outputs and nucleus sampling ($T{=}0.7$, $p{=}0.9$) for
the unwatermarked baseline. Primary model: Llama-3.2-3B-Instruct;
secondary models for cross-model runs: Mistral-7B, GPT2-XL. Secret
key \texttt{curated\_wiki\_dataset\_2024}, seed 42, max 200 tokens.

We report six experiments. E1--E3 validate the calibration formula
and the quality/robustness theorems directly. E4 is the critical
test of the low-entropy-robustness hypothesis. E5--E6 validate
self-healing and HDD.

\paragraph{E1: Calibration.} We verify that $S^\star(n, \rho, \alpha)$
predicts empirical TPR and FPR. Configurations span
$S \in \{2, 3, 4, 5, 7, 9\}$, entropy thresholds $\tau_H \in
\{1.5, 2.0, 2.5, 3.0, 3.5\}$, text lengths $n \in \{200, 500\}$.
The prediction is that $S \geq S^\star(n, \rho, \alpha)$ achieves
$\mathrm{TPR} \approx 1$ and $\mathrm{FPR} \approx \alpha$; $S < S^\star$
misses the target.

\paragraph{E2: Quality.} Matched-detection comparison between
entropy-gated configurations on the calibration curve and a non-gated
reference ($\rho = 1$, $S = 7$). At equal detection power (equal $z$),
does gated MCL have lower perplexity than the non-gated reference?
\cref{thm:quality} predicts yes for sufficiently long texts.

\paragraph{E3: Robustness.} Seven attacks:
(i) random word replacement, $\delta \in \{0.1, 0.2, 0.3, 0.4, 0.5\}$;
(ii) WordNet synonym substitution; (iii) DIPPER paraphrase;
(iv) back-translation (en$\to$fr$\to$en, en$\to$zh$\to$en);
(v) homoglyph perturbation; (vi) word insertion; (vii) word deletion.
Primary metric: TPR at $p < 0.01$ after attack. \cref{thm:robustness}
predicts $\delta^\star \approx 0.293$ universal for (i); (iii)--(iv)
are the real-world rewrite tests that load on E4.

\paragraph{E4: Entropy--survival (critical).} For 500 watermarked
texts, record the per-position entropy $H_i$ during generation.
Paraphrase each text with DIPPER, GPT-rewrite, and back-translation.
Align original and rewritten tokens via edit distance. For each
entropy bin, compute the empirical \emph{token survival rate}
$\Pr[t'_i = t_i \mid H_i]$. The hypothesis is that survival rate
decreases monotonically with $H_i$ and exceeds $0.6$ for $H_i < 1.0$.
\textbf{Falsification would invalidate the low-entropy-robustness
claim} and force a revision of \S\ref{sec:discussion}.

\paragraph{E5: Self-healing.} For
$S \in \{2, 3, 5, 7\}$ and $\delta \in \{0.1, 0.2, 0.3, 0.4, 0.5\}$,
compare the observed $\phi$ after random replacement to the closed
form in \cref{thm:robustness}, and verify $\delta^\star \approx 0.293$
is universal across $S$.

\paragraph{E6: HDD vs.\ model-free detection.} ROC comparison of
three detectors: (a) oracle-entropy HDD, (b) proxy-entropy HDD using
GPT-2 as a cheap verifier for Llama-3.2, and (c) model-free
detection. \cref{thm:hdd} predicts (a) $>$ (c); the open empirical
question is how much (b) concedes relative to (a).

\paragraph{Reference point.} A non-gated configuration at $S{=}7$,
$\rho{=}1$ reaches $100\%$ TPR, $0\%$ FPR, perplexity $4.20$;
unwatermarked baseline perplexity is $\sim 3.0$.
\cref{tab:baseline} in \cref{app:results} gives the full sweep.

% Current status: E1/E2/E5 partial results in progress;
% E3/E4/E6 outstanding.
